\subsection{实验结果与分析}

\subsection*{（一）实验数据}

实验测得的数据如表~\ref{tab:data} 所示。

\begin{table}[htbp]
    \centering
    \caption{实验数据记录（示例）}
    \label{tab:data}
    \begin{tabular}{cccc}
        \toprule
        \textbf{序号} & \textbf{输入} & \textbf{预期输出} & \textbf{实际输出} \\
        \midrule
        1 & 0x01 & 0x02 & 0x02 \\
        2 & 0x03 & 0x06 & 0x06 \\
        3 & 0x07 & 0x0E & 0x0E \\
        \bottomrule
    \end{tabular}
\end{table}

\subsection*{（二）结果图示}

单张图片的插入方式如图~\ref{fig:example} 所示（此处以校徽图片为例，实际使用时替换为实验截图或波形图）。

\begin{figure}[htbp]
    \centering
    \includegraphics[width=0.35\textwidth]{zjulogo.png}
    \caption{示例图片}
    \label{fig:example}
\end{figure}

多张子图的排布方式如图~\ref{fig:sub} 所示，其中子图~\subref{fig:sub-a} 与子图~\subref{fig:sub-b} 分别展示两组对比结果。

\begin{figure}[htbp]
    \centering
    \begin{subfigure}[b]{0.4\textwidth}
        \centering
        \includegraphics[width=0.6\textwidth]{zjulogo.png}
        \caption{子图一}
        \label{fig:sub-a}
    \end{subfigure}
    \hfill
    \begin{subfigure}[b]{0.4\textwidth}
        \centering
        \includegraphics[width=0.6\textwidth]{zjulogo.png}
        \caption{子图二}
        \label{fig:sub-b}
    \end{subfigure}
    \caption{子图示例}
    \label{fig:sub}
\end{figure}

\subsection*{（三）结果分析}

结合表~\ref{tab:data} 与图~\ref{fig:example}，对实验结果进行分析：实际输出与预期输出一致，验证了设计的正确性。分析方法可参考文献~\cite{smith2020analysis}。